% ACKNOWLEDGEMENTS
\chapter*{Acknowledgements}

The story of this thesis begins a year ago, in March 2024, landing in Kyoto during the Sakura season (cherry blossom). The blooming of my research has lasted for a year since that time, and it has been the period of my academic studies in which I started to love what I do. I want to thank Professor Yoshida Kazuya that accept me as a research student at the Space Robotics Lab, giving me the possibility to work on something unique in the world. I would like to thank my advisor, Professor Luca Bascetta, for his support in critical moments of this research.
My work in Japan would not be possible without the help of Associate Professors Shreya Santra and Kentaro Uno, their continuous feedback and their kind words of appreciation gave me the motivation to keep going. I also want to thank the members of the SRL who support my vision there, especially Takada and Naoki, who helped me during my stay and continue to exchange ideas. I would love to work together again, pushing the boundaries of robotics research in exploring the unknown. 

Adesso i ringraziamenti entrano nella fase emotiva...Voglio dedicare questa tesi a mia nonna Renate, che mi ha trasmesso la sua determinazione e un filtro da artista sul mondo. Ai miei genitori (mamma, papà e zia Luisa – sì, io ne ho tre), che sono un faro per me e mi hanno insegnato che, ogni tanto, dei sacrifici sono necessari (come mi ha dimostrato questa tesi), ma che le soddisfazioni eclissano ogni fatica. La pazzia e l’ostinatezza di mio padre, la precisione di mia zia, la curiosità e l’instancabilità di mia madre sono state fondamentali, e mi accompagneranno sempre. 

Cecilia mi ha permesso di fare una seconda laurea in medicina, vivendo tutte le ansie di un esame di semeiotica o eziopatologia. Lei è stata la mia più grande fan, come io sono il suo, e ha reso Milano una bella città in cui ritornare dopo ogni festività passata giù. Grazie a lei questi anni sono stati pieni di colori, sorrisi, ed esperienze, anche quando il cielo era grigio. Questo primo passo verso una nuova fase della mia vita lo dedico anche a te. 

La cuginanza e gli zii andrebbero ringraziati uno ad uno: mi hanno fatto sentire a casa. Tra pranzi domenicali in famiglia, birrette,e serate in bici vagando per milano, senza di loro questi cinque anni non sarebbero stati gli stessi. Ovviamente va ringraziata anche la rimanente parte della cuginanza e i miei amici in Sicilia, per avermi dato un motivo per sopravvivere fino all’estate ogni anno.
Trasferirsi a Milano mi ha avvicinato di più alla parte della famiglia "emigrata", riducendo finalmente la distanza fisica che ci ha sempre separato. Grazie Bea, che mi hai reso una persona più critica, incazz**** verso le inguistizie e con un occhio sempre attento a ciò che mi succede attorno. Grazie a los Santos e Iago, fratelli e compagni di vita, tra pazzia, nosense ed esperienze indimenticabili. Cami e Ludo, tra iperattività e "ciaeee", Sere, Giuli, Iaffio e Vivi, grazie a tutti. Grazie anche ad un Siculo-Arabo-Ispanico, e tutti i miei amici che da ormai più di 5 anni aspettano ogni mio ritorno a casa.  

Le persone incontrate durante questi anni al Politecnico che ho ancora accanto sono ormai una parte integrante della mia vita, e non intendo chiudere questi rapporti con le conclusioni di questa tesi.
Pazzi di biorobot, fanatici di treni, gymbro, gattini dell'Agorà, amanti delle matrici, e chi più ne ha più ne metta, tutto questo non sarebbe stato lo stesso senza di voi. 

Se tu, lettore, stai per leggere questa tesi, ti auguro buona fortuna e ti ringrazio. Questo lavoro è stato una soddisfazione, seguito da mille problemi, e poi di nuovo soddisfazioni. Sono fiero dei risultati ottenuti e spero che il lavoro svolto possa ispirarti, come ha fatto con me.

